\documentclass{amsart}

\usepackage{keytheorems}
\usepackage{amsthm}
\usepackage{color}

\newkeytheorem{problem}
\newtheorem*{proposition}{Proposition}
\newcommand{\todo}{\textcolor{red}{TO-DO.} }
\newcommand{\temp}{\textcolor{red}{TEMPORARY.} }

\newcommand{\set}[1]{\left\{#1\right\}}
\newcommand{\powset}{{\cal P}}
\newcommand{\img}{_{*}}
\newcommand{\pimg}{^{*}}
\newcommand{\limn}{\displaystyle \lim_{n \rightarrow \infty}}

\newcommand{\R}{\mathbb{R}}
\newcommand{\Z}{\mathbb{Z}}
\newcommand{\C}{\mathbb{C}}
\newcommand{\N}{\mathbb{N}}
\newcommand{\Q}{\mathbb{Q}}

\numberwithin{equation}{section}

\begin{document}

	\title[Math 251W: Homework \#12]{
		Math 251W: Foundations of Advanced Mathematics \\
		\textmd{Homework \#12: Portfolio Problems \S 6.4 \& 7.8}}
	\author{Carmen Beltran, Brendan Butler, Teddy Wachtler}
	\date{\today}

	\maketitle

	\begin{problem}[manual-num=63]
		\begin{proposition}
			Let $r_1, r_2, r_3, \dots$ be the sequence defined by $r_1 = 1$ and $r_{n + 1} = 4r_n + 7$ for all natural numbers $n$.
			Prove that $r_n = \frac{1}{3} (10 \cdot 4^{n-1} - 7)$ for all natural numbers $n$.
		\end{proposition}

		\begin{proof}
			\todo
		\end{proof}
	\end{problem}

	\begin{problem}[manual-num=69]
		\begin{proposition}
			Use the definition of a limit to prove $\limn \frac{n^2}{5-n^2} = -1$.
		\end{proposition}

		\begin{proof}
			\todo
		\end{proof}
	\end{problem}

	\begin{problem}[manual-num=74]
		\begin{proposition}
			Use strong induction to prove that every natural number can be written as a sum of distinct powers of 2.
		\end{proposition}

		\begin{proof}
			\todo
		\end{proof}
	\end{problem}

\end{document}